\section*{Hoofdstuk \ref{ch:g11:genes-and-disease} --- Genetische variatie en ziekte}

\begin{solution}{exo:g11:genes-and-disease:1}
\omterm{def:g11:genes-and-disease:genetic}{Recessief}: de ziekte verschijnt alleen met twee kopieën van het \omterm{def:g10:universal-dna:gene}{allel};
\omterm{def:g11:genes-and-disease:genetic}{dominant}: één kopie volstaat. Een drager heeft één \omterm{def:g11:genes-and-disease:genetic}{recessief} gemuteerd
\omterm{def:g10:universal-dna:gene}{allel} en is gezond.
\end{solution}

\begin{solution}{exo:g11:genes-and-disease:2}
Aangedaan $1/4$, drager $1/2$, geen van beide $1/4$.
\end{solution}

\begin{solution}{exo:g11:genes-and-disease:3}
I-1 $Aa$, I-2 $Aa$ (ouders van een aangedaan kind), II-2 $aa$, III-2
$aa$.
\end{solution}

\begin{solution}{exo:g11:genes-and-disease:4}
Een man heeft één \omterm{def:g11:cell-cycle-mitosis:chromosome}{X-chromosoom}: een \omterm{def:g11:genes-and-disease:genetic}{recessief} \omterm{def:g10:universal-dna:gene}{allel} daarop komt tot
uiting, want er is geen tweede kopie om het te maskeren. Een vrouw heeft
het \omterm{def:g10:universal-dna:gene}{allel} op allebei haar \omterm{def:g10:universal-dna:information}{X-chromosomen} nodig.
\end{solution}

\begin{solution}{exo:g11:genes-and-disease:5}
Een ziekte waarvan het risico van vele \omterm{def:g10:universal-dna:gene}{allelen} met een klein effect en van
de omgeving afhangt: diabetes type 2, hartziekte (en ook astma, de meeste
vormen van kanker).
\end{solution}

\begin{solution}{exo:g11:genes-and-disease:6}
$aa \times AA$: alle kinderen $Aa$, gezonde dragers. Een dragend kind met
een drager: $1/4$ aangedaan, $1/2$ drager, $1/4$ niet-drager.
\end{solution}

\begin{solution}{exo:g11:genes-and-disease:7}
II-1 is een gezond kind van twee dragers: drager met kans $2/3$. Risico
met een man uit de bevolking: $2/3 \times 1/25 \times 1/4
= 1/150$.
\end{solution}

\begin{solution}{exo:g11:genes-and-disease:8}
$1/2$ voor elk kind. Geen van de vier: $(1/2)^4 = 1/16$.
\end{solution}

\begin{solution}{exo:g11:genes-and-disease:9}
$X^hY \times X^HX^H$: alle dochters $X^HX^h$, gezonde draagsters; alle
zonen $X^HY$, gezond. Een dragende dochter met een gezonde man: de helft
van haar zonen heeft hemofilie, de helft van haar dochters is draagster.
\end{solution}

\begin{solution}{exo:g11:genes-and-disease:10}
Het ene \omterm{def:g10:universal-dna:gene}{allel} geeft 400 en 200; het andere, met de extra deletie, geeft
300 en 200. Banden op 400, 300 en 200 --- en geen 600.
\end{solution}

\begin{solution}{exo:g11:genes-and-disease:11}
De frequentie van een \omterm{def:g10:universal-dna:gene}{allel} in een bevolking hangt van haar geschiedenis
af en niet alleen van het mutatietempo: het Europese \omterm{def:g10:universal-dna:gene}{allel} ontstond lang
geleden en verbreidde zich, misschien omdat dragers enig voordeel hadden
(weerstand tegen een diarreeziekte wordt vermoed); in Oost-Azië gebeurde
dat niet.
\end{solution}

\begin{solution}{exo:g11:genes-and-disease:12}
Elk kind is een onafhankelijke trekking: $1/4$ voor het tweede. Ten minste
één van de volgende twee: $1 - (3/4)^2 = 7/16$. De \omterm{def:g10:universal-dna:gene}{allelen} worden bij elke
verwekking opnieuw gedeeld; een eerder aangedaan kind verlaagt het risico
van het volgende niet.
\end{solution}

\begin{solution}{exo:g11:genes-and-disease:13}
Vóór de test is zij drager met kans $1/25$; de test mist 10\% van de
dragers, dus zijn van 1000 vrouwen er 40 drager en testen er 4 negatief,
terwijl 960 niet-dragers negatief testen: de kans om drager te zijn is
$4/964 \approx 1/240$. Risico met een bekende drager: $1/240 \times 1/4
\approx 1/960$.
\end{solution}

\begin{solution}{exo:g11:genes-and-disease:14}
Van veel naar weinig \omterm{def:g10:universal-dna:gene}{allelen}, zittend: 34\% naar 9\%; van zittend naar
actief, met veel \omterm{def:g10:universal-dna:gene}{allelen}: 34\% naar 12\%. De verandering die je wél kunt
maken, geeft vrijwel dezelfde daling als die welke je niet kunt maken. Een
risicoscore is nuttig om te bepalen wie het meest wint bij een andere
leefwijze, niet als vonnis.
\end{solution}

\begin{solution}{exo:g11:genes-and-disease:15}
De test biedt zekerheid --- verlost zijn van een schaduw van vijftig
procent bij een negatieve uitslag, en de mogelijkheid om een leven en een
gezin te plannen bij een positieve --- en kost, bij een positieve
uitslag, de wetenschap van een ongeneeslijke ziekte tientallen jaren van
tevoren, met gevolgen voor de stemming, verzekeringen en verwanten die het
\omterm{def:g10:universal-dna:gene}{gen} delen. Er volgt geen behandeling uit de uitslag, dus de enige reden om
het te weten is die van de persoon zelf; en daarom liggen de keuze en het
tijdstip ervan bij die persoon.
\end{solution}

\begin{solution}{pb:g11:genes-and-disease:1}
\textbf{1.} \omterm{def:g11:genes-and-disease:genetic}{Recessief}: Zoé is aangedaan en haar beide ouders zijn gezond,
dus draagt elk van hen één verborgen gemuteerd \omterm{def:g10:universal-dna:gene}{allel}.

\textbf{2.} Léa $Aa$, Tom $Aa$, Zoé $aa$.

\textbf{3.} $AA$ met kans $1/3$, $Aa$ met $2/3$ (hij is gezond, dus $aa$
valt af).

\textbf{4.} Tom is drager, dus is ten minste een van zijn ouders dat ook;
omdat er in die generatie niemand aangedaan is, kreeg Ana het \omterm{def:g10:universal-dna:gene}{allel} van
die ouder met kans $1/2$: $1/2$ (met de andere ouder als niet-drager
genomen).

\textbf{5.} $1/25$, het cijfer van de bevolking.

\textbf{6.} $1/2 \times 1/25 \times 1/4 = 1/200$.

\textbf{7.} Aangedaan $1/4$; drager $1/2$.

\textbf{8.} $2/3 \times 1/25 \times 1/4 = 1/150$.

\textbf{9.} Zoé is $aa$; haar partner is $Aa$ met kans $1/25$. Aangedaan:
$1/25 \times 1/2 = 1/50$. Drager: elk kind krijgt $a$ van Zoé, dus is het
drager met kans $24/25 \times 1 + 1/25 \times
1/2 \approx 0.98$.

\textbf{10.} Lou is drager met kans $2/3$; de neef of nicht kreeg $a$ van
Ana ($1/2$ drager) met kans $1/2$, en is dus drager met kans $1/4$ (Karim
buiten beschouwing gelaten): $2/3 \times 1/4
\times 1/4 = 1/24$, vier keer hoger dan bij Max, omdat de neef of nicht
uit dezelfde familie komt en veel eerder hetzelfde \omterm{def:g10:universal-dna:gene}{allel} draagt dan een
buitenstaander.

\textbf{11.} Zoé: alleen 600 en 300. Tom: 900, 600 en 300. Niet-drager:
alleen 900.

\textbf{12.} $AA$: hij is geen drager; het risico voor zijn kind daalt van
$1/150$ naar nul (behoudens een nieuwe \omterm{def:g11:mutations:mutation}{mutatie}).

\textbf{13.} Ana is $Aa$: risico $1 \times 1/25 \times 1/4 = 1/100$.

\textbf{14.} Van 1000 mensen zoals Karim zijn er 40 drager, van wie er 4
negatief testen, en testen 960 niet-dragers negatief: $4/964 \approx
1/240$. Risico: $1 \times 1/240 \times 1/4 \approx 1/960$.

\textbf{15.} $Aa$: een gezonde drager, net als zijn moeder.

\textbf{16.} $(1/25)^2 = 1/625$ van de paren bestaat uit twee dragers; een
kwart van hun kinderen is aangedaan: $1/2500$.

\textbf{17.} Precies de waargenomen frequentie: de dragersfrequentie en de
recessieve regel verklaren haar.

\textbf{18.} Bij een dragersfrequentie van $1/25$ en een frequentie van
aangedanen van $1/2500$ overtreffen de gemuteerde \omterm{def:g10:universal-dna:gene}{allelen} bij gezonde
dragers die bij aangedane personen met $2 \times 1/25$ tegen $2 \times
1/2500$, honderd tegen één. De dood van de aangedanen verwijdert per
generatie maar 1\% van de kopieën; de rest wordt onzichtbaar door dragers
doorgegeven.

\textbf{19.} De frequentie van het \omterm{def:g10:universal-dna:gene}{allel} zal stijgen, maar heel traag: de
1\% van de kopieën die vroeger elke generatie verloren ging, blijft nu
behouden, dus is de verandering van de orde van 1\% van de frequentie per
generatie.

\textbf{20.} Twee dragers lopen één kans op vier op een aangedaan kind;
een gezonde broer of zus van een aangedaan kind is twee keer op drie
drager; de gel maakte van kansen \omterm{def:g11:enzymes-and-phenotype:phenotype}{genotypen} --- Max vrijgepleit, Ana
bevestigd, Karim gerustgesteld maar niet vrijgepleit --- en wees het
ongeboren kind aan als gezonde drager.
\end{solution}
